% Generated draft: requires mathematical and source review.
\documentclass[11pt]{article}
\usepackage[margin=25mm]{geometry}
\usepackage{fontspec}
\setmainfont{Latin Modern Roman}
\usepackage{amsmath,amssymb,mathtools,unicode-math}
\setmathfont{Latin Modern Math}
\usepackage{enumitem,needspace,xurl,hyperref}
\setlength{\parindent}{0pt}
\setlength{\parskip}{5pt}
\setlength{\emergencystretch}{3em}
\newcommand{\N}{\mathbb{N}}
\newcommand{\Z}{\mathbb{Z}}
\newcommand{\Q}{\mathbb{Q}}
\newcommand{\R}{\mathbb{R}}
\newcommand{\C}{\mathbb{C}}
\newcommand{\F}{\mathbb{F}}
\newcommand{\A}{\mathbb{A}}
\newcommand{\PP}{\mathbb{P}}
\newcommand{\E}{\mathbb{E}}
\newcommand{\eps}{\varepsilon}
\newcommand{\dd}{\,\mathrm{d}}
\begin{document}
\section*{UnsolvedMath ID: 30005653}
\textbf{Source problem code:} catalogue entry
\subsection{Definitions and mathematical statement}
In the usual module-theoretic setting, a stable module category has modules as objects and morphisms given by module homomorphisms modulo those factoring through projective modules. When a group algebra has a block decomposition, its principal block is the block containing the trivial module. An elementary abelian group is a finite direct product of cyclic groups of one prime order. These descriptions explain the terminology; the supplied excerpt does not specify the groups, coefficient fields, characteristics, or size conventions for the stable module categories under consideration.

\textbf{Main question (Question 12 in the supplied quotation).} Do the spectral schemes occurring in the source's preceding equivalences admit natural spectral moduli-theoretic interpretations? The supplied cleaned statement identifies these as spectral schemes associated with principal blocks of stable module categories and asks for such interpretations beyond the elementary abelian case.

A moduli-theoretic interpretation should specify the mathematical objects being parametrized, their families over suitable test objects, their equivalences, and the correspondence with the associated spectral scheme. However, the preceding equivalences and the construction of these spectral schemes are absent from the supplied record. The meaning of ``spectral scheme,'' the category of test objects, and the required notion of naturality are also unspecified. Consequently, this excerpt does not determine a fully specified representability statement or its precise range of quantification. These missing data must be recovered before a precise mathematical formulation can be completed.

\textbf{Supplied qualification.} The unidentified speaker states that they can presently give an affirmative answer only for elementary abelian groups. This is a source-relative report of available constructions, not a claim that other groups cannot admit such interpretations.

\textbf{Adjacent question retained from the source (Question 11).} Can the Picard and Brauer groups of the principal blocks for stable module categories be completely computed by a descent spectral sequence based on their associated spectral schemes? Conventionally, a Picard group concerns isomorphism classes of tensor-invertible objects, and a Brauer group concerns Morita-equivalence classes of suitable Azumaya algebra objects. The exact categorical conventions, descent construction, and convergence conditions intended here are not supplied. This adjacent computational question is preserved separately; its equivalence to Question 12 is not asserted.
\subsection{Short English statement}
The main question asks whether spectral schemes attached to principal blocks of stable module categories naturally classify mathematical objects and their families, beyond the elementary abelian case. The omitted preceding equivalences and unspecified moduli problem prevent a complete precise formulation. The source also includes a separate question about computing Picard and Brauer groups by descent.
\subsection{Sources}
\begin{itemize}
\item Supplied local UnsolvedMath registry record, permanent ID \(30005653\), source problem code \texttt{catalogue entry}, title ``Spectral Moduli Interpretations for Stable Module Blocks.'' The record supplies proposed year 2023 and no named proposer; the first-person attribution in Question 12 remains unidentified.
\item Supplied source citation: \emph{Tensor-Triangular Geometry and Interactions (2024)}, Oberwolfach Reports, DOI 10.4171/owr/2023/40, pp.~2303--2358. \url{https://doi.org/10.4171/owr/2023/40}. The quotation contains Questions 12 and 11 and refers to preceding equivalences not included in the record.
\item Supplied literature reference: Benson--Iyengar--Krause, ``stratification for finite group schemes'' (label supplied by the registry). \url{https://arxiv.org/abs/1510.06773}.
\item Literature assessment dated exactly \texttt{2026-08-22}: the supplied curation labels the problem ``partially solved,'' reports natural spectral interpretations for elementary abelian groups and related support constructions, and states that the general moduli-theoretic interpretation question remains unanswered in the report. This is supplied dated curation, not an independently verified literature assessment; its time-sensitive status and cited claims require independent verification.
\item The registry labels the statement \texttt{corrected\_verified} and reports direct checking of the official report. That verification claim is supplied provenance and has not been independently checked for this draft.
\end{itemize}
\end{document}
